\section{The target-directed progressive calculus}
\label{target:section}

The arithmetic terms and the progressive interpretation of atomic closures
are inherited.  The new change concerns elimination: a closure instruction
supplies only its desired target term.  A syntactic successor target asks for
its predecessor instance; every other atomic target asks for the zero
instance.  An iterator with zero counter is a certified way to request the
second mode even when the desired result is a successor.  This removes the
separate numerical source parameter and the addition term used by the
\texttt{uniform\_progress} proposal.

The construction below has seven primitive inference schemas and the full,
unrestricted equation/$Q$/$D$ formula grammar.  Its theorem is exact
consequence equivalence with that earlier progressive calculus, including
open sequents and every formula of their common language.  It is a further
presentation refinement, not a new arithmetic foundation or a claim of
absolute minimality.

\subsection{Complete primitive specification}
\label{target:specification}

There is one numerical sort.  Terms and formulas are finite syntax trees:
\[
 t ::= x\mid 0\mid S(t)\mid I(n,a;u.h),
 \qquad
 A ::= (s=t)\mid Q_y A\mid D_{x,y}A.
\]
Only $u$ in $h$ is bound by the iterator.  The closure $Q$ binds $y$; $D$
binds the distinct variables $x,y$.  All substitution is capture-avoiding.
Terms are not silently normalized before a rule is selected.  In particular,
$I(0,Sr;u.0)$ is an iterator-headed term, although its value is $Sr$.
The contradiction is $0=S0$.

Call a matrix \emph{atomic} when its outermost constructor is equality,
and \emph{compound} when its root is $Q$ or $D$.  Set
\[
 \eta_B(z)=
 \begin{cases}
 z,&B\text{ atomic},\\
 0,&B\text{ compound},
 \end{cases}
 \qquad
 \delta_B(t)=
 \begin{cases}
 r,&B\text{ atomic and }t\text{ is literally }Sr,\\
 0,&\text{otherwise}.
 \end{cases}
 \tag{Target dispatch}\label{target:dispatch}
\]
These are metasyntactic definitions, not new numerical functions or tests
for whether an arbitrary number is zero or a successor.

Contexts are finite reusable lists of formulas.  The calculus $\mathcal T$
has precisely the following seven primitive inference schemas:
\[
 \frac{A\in\Gamma}{\Gamma\vdash_{\mathcal T}A}
 \quad\mathrm{Hyp},
\]
\[
 \frac{}{\Gamma\vdash_{\mathcal T}I(0,a;u.h)=a}
 \quad\mathrm{I0},
 \qquad
 \frac{}{\Gamma\vdash_{\mathcal T}
 I(Sn,a;u.h)=h[I(n,a;u.h)/u]}
 \quad\mathrm{IS},
\]
\begin{equation}
 \frac{\Gamma\vdash_{\mathcal T}s=t
       \qquad\Gamma\vdash_{\mathcal T}p(s)=q(s)}
      {\Gamma\vdash_{\mathcal T}p(t)=q(t)}
 \quad\mathrm{Tr}_{=},
 \label{target:transport-rule}
\end{equation}
\begin{equation}
 \frac{\Gamma\vdash_{\mathcal T}Q_yB(y)
       \qquad\Gamma\vdash_{\mathcal T}B(\delta_B(t))}
      {\Gamma\vdash_{\mathcal T}B(t)}
 \quad\mathrm{QE}_{\rm target},
 \label{target:elimination-rule}
\end{equation}
\begin{equation}
 \frac{\Gamma\vdash_{\mathcal T}D_{x,y}B(x,y)
       \qquad\Gamma,Q_yB(b,y)\vdash_{\mathcal T}A}
      {\Gamma\vdash_{\mathcal T}A}
 \quad\mathrm{DE},
 \label{target:de-rule}
\end{equation}
\begin{equation}
 \frac{\Gamma,B(w,\eta_B(z))\vdash_{\mathcal T}B(w,Sz)}
      {\Gamma\vdash_{\mathcal T}D_{x,y}B(x,y)}
 \quad\mathrm{DI}_{\rm root}.
 \label{target:di-rule}
\end{equation}
Here $p(v),q(v)$ in \eqref{target:transport-rule} are numerical term
templates.  In DE, $b$ is fresh for $\Gamma,A$ and the fixed parameters of
$B$.  In DI, $z$ is fresh for $\Gamma,w$ and the fixed parameters of $B$.
The displayed local assumptions are discharged.  Weakening, contraction,
exchange, cut, and numerical substitution are admissible operations, not
additional primitive instructions.

There are no other primitive logical axioms, ordinary logical connectives,
quantifier rules, successor axioms, equality axioms, or induction schemas.
The literal test in \eqref{target:dispatch} is part of the complete
specification; replacing it by a test of the value of $t$ would define a
different system.

\subsection{The two closure interfaces and induction}
\label{target:bootstrap}

Weakening, exchange, contraction, fresh renaming, and assumption cut follow
by the usual induction on finite derivations.  For cut, replace the
appropriate hypothesis leaves and rename eigenparameters first.  Numerical
substitution needs a separate check and is proved below, after the elementary
closure construction it uses.

Put
\begin{equation}
 W(t):=I(0,t;u.0),\qquad 1:=S0.
 \label{target:wrapper}
\end{equation}
The state variable is fresh and unused.  The equation $W(t)=t$ is an I0
instance, while $W(t)$ is never syntactically a successor.

\begin{lemma}[Equality and derived closure introduction]
\label{target:eq-qi}
Reflexivity, symmetry, transitivity, and numerical congruence are derivable
in $\mathcal T$.  In addition,
\[
 \frac{\Gamma,B(\eta_B(z))\vdash_{\mathcal T}B(Sz)}
      {\Gamma\vdash_{\mathcal T}Q_yB(y)}
 \quad\mathrm{QI}_{\rm root}
\]
is admissible, with $z$ fresh for the retained context and fixed parameters.
\end{lemma}
\begin{proof}
Use $W(a)=a$ as both premises of equational transport with equation template
$v=a$ to obtain $a=a$.  From $s=t$ and $s=s$, transport with template
$v=s$ gives $t=s$.  The usual equation templates then give transitivity and
numerical congruence; no transport through a compound formula is used.

For QI, apply DI at witness $0$ to $D_{r,y}B(y)$, with $r$ absent from
$B$.  Eliminate this result by DE.  Its opened assumption is $Q_yB(y)$,
independent of the fresh witness parameter, so Hyp supplies the desired
conclusion of the DE minor.
\end{proof}

\begin{lemma}[One step and arbitrary-target closure]
\label{target:closure-interfaces}
For every equation $E(y)$ and every term $t$, the following rules are
admissible under arbitrary retained contexts:
\[
 Q_yE(y),\ E(t)\ \vdash_{\mathcal T}\ E(St),
 \qquad
 Q_yE(y),\ E(0)\ \vdash_{\mathcal T}\ E(t).
 \tag{Atomic interfaces}\label{target:atomic-interfaces}
\]
For every compound matrix $B$, ordinary closure elimination is primitive:
\[
 Q_yB(y),\ B(0)\ \vdash_{\mathcal T}\ B(t).
 \tag{Compound interface}\label{target:compound-interface}
\]
\end{lemma}
\begin{proof}
For one-step progress, choose target $St$ in
\eqref{target:elimination-rule}; its selected source is exactly $t$.

For arbitrary-target atomic closure, first choose target $W(t)$.  Since it
is iterator-headed, its selected source is $0$.  QE gives $E(W(t))$ from
$Q_yE$ and $E(0)$.  The I0 equation $W(t)=t$ and primitive equational
transport then give $E(t)$.  This derivation uses no numerical substitution
metatheorem and no compound transport.

For compound $B$, the definition of $\delta_B$ is identically $0$, so the
displayed rule is a direct QE instance.
\end{proof}

\begin{corollary}[Contextual equation induction]
\label{target:equation-induction}
For every literal equation $E(y)$,
\[
 \frac{\Gamma\vdash_{\mathcal T}E(0)
       \qquad\Gamma,E(z)\vdash_{\mathcal T}E(Sz)}
      {\Gamma\vdash_{\mathcal T}E(t)}
\]
is admissible, where $z$ is fresh for the retained context and fixed
parameters.
\end{corollary}
\begin{proof}
The step premise gives $Q_yE(y)$ by QI.  Apply the arbitrary-target atomic
closure from Lemma~\ref{target:closure-interfaces} to the supplied base.
\end{proof}

This is the complete bootstrap.  It requires neither a predecessor, a
numerical case split, an identity for addition, nor an assumption that a
syntactically nonsuccessor term denotes zero.  The wrapper in
\eqref{target:wrapper} is a constant-size term constructor in the proof
translation; its numerical value need not be computed by the checker.

\begin{lemma}[Capture-free numerical substitution]
\label{target:substitution}
If $\Gamma\vdash_{\mathcal T}A$, then
$\Gamma\sigma\vdash_{\mathcal T}A\sigma$ for every capture-free numerical
substitution $\sigma$.
\end{lemma}
\begin{proof}
Induct on the derivation, freshly renaming discharged eigenparameters.
All rules other than QE commute with substitution in their displayed forms:
in particular, the outer formula constructor, and hence $\eta_B$, is
preserved.

Consider QE with matrix $B$ and target $t$.  If $B$ is compound, both before
and after substitution the source is $0$.  If $B$ is an equation and
$t=Sr$, then $t\sigma=S(r\sigma)$ and the selected source remains
$r\sigma$.  If $B$ is an equation and neither $t$ nor $t\sigma$ is a
syntactic successor, both selected sources are $0$.  These cases translate
to primitive QE instances.

The remaining case is an atomic matrix with a nonsuccessor target $t$
whose substitution becomes a successor.  The translated premises are
$Q_y(B\sigma)(y)$ and $(B\sigma)(0)$, whereas a primitive QE at $t\sigma$
would request its predecessor instance.  Apply the already proved
arbitrary-target atomic closure instead.  It yields $(B\sigma)(t\sigma)$
from exactly those translated premises.  Lemma~\ref{target:closure-interfaces}
was established without the present substitution theorem, so this repair
does not introduce a circular dependency.
\end{proof}

\subsection{Transport through the whole formula language}
\label{target:full-transport-section}

\begin{lemma}[Full transport and homogeneous binder congruence]
\label{target:full-transport}
For every native formula template $F(v)$,
\[
 s=t,\ F(s)\ \vdash_{\mathcal T}\ F(t).
\]
If matrices $B,C$ are uniformly interderivable and have the same root class
(both equations or both compound), their $Q$ closures are interderivable,
as are their $D$ closures.
\end{lemma}
\begin{proof}
For transport, induct on the formula template.  An equation uses primitive
equational transport.  For $Q_yE(v,y)$ with an equation matrix, introduce
the target closure by atomic QI: move its temporary $E(t,z)$ back to
$E(s,z)$, use one-step elimination for the retained source closure, and
move $E(s,Sz)$ forward to $E(t,Sz)$.  Both moves are equation transports.

For $Q_yB(v,y)$ with a compound matrix, introduce the target by compound
QI.  The induction hypothesis moves its local base $B(t,0)$ to $B(s,0)$;
compound QE at $Sz$ gives $B(s,Sz)$; the induction hypothesis moves this
back to $B(t,Sz)$.  These transports concern the strictly smaller matrix.

For $D_{x,y}B(v,x,y)$, eliminate the source at a fresh witness $b$ and
introduce the target at that same witness.  Use the induction hypothesis
to move the local introduction premise from parameter $t$ to $s$.  Apply
one-step elimination if the matrix is atomic and ordinary compound closure
elimination otherwise.  Move the successor conclusion back by the induction
hypothesis, and close DI and DE.  Every transport used inside this
construction concerns a strictly smaller formula.

For homogeneous congruence, the same constructions use the two supplied
matrix conversions in place of numerical transport.  The root-class
condition ensures that the two DI or QI premises use the same source
$\eta(z)$.  In the $D$ case, eliminate the source binder and reuse its fresh
witness for the target binder.
\end{proof}

The qualification ``homogeneous'' matters.  An equation can be equivalent
to a compound formula without its progressive closure being equivalent to
the compound formula's conditional-universal closure.  No unrestricted
congruence across these two root classes is asserted.

\subsection{Exact comparison with the inherited orbit calculus}
\label{target:orbit-equivalence-section}

Let $\mathcal R$ be precisely the seven-rule calculus in the incoming
\texttt{uniform\_progress} report.\cite{uniform}  It has the same terms, formulas, Hyp,
I0, IS, equational transport, DE, and root-directed DI as $\mathcal T$.
Only its closure elimination differs.  Write
\[
 A(a,n):=I(a,n;u.Su).
\]
Its elimination rule is
\begin{equation}
 \frac{Q_yB(y)\qquad B(\eta_B(a))}
      {B(A(\eta_B(a),n))}.
 \tag{Orbit QE}\label{target:old-qe}
\end{equation}
The formulas and terms in this comparison have their literal common
syntax; the translations between $\mathcal T$ and $\mathcal R$ are
identity maps on them, with changes only to their proof certificates.

\paragraph{Dependencies on the $\mathcal R$ side.}
The inherited report first proves contextual equation induction inside
$\mathcal R$: at source $0$, \eqref{target:old-qe} gives
$E(I(0,t;u.Su))$, and I0 plus equational transport gives $E(t)$.
It then proves the elementary addition identities by equation induction,
extracts $Q E,E(a)\vdash E(Sa)$ using $A(a,1)=Sa$, and derives full
formula transport by structural induction.  In particular, its compound
transport is available before the comparison below.  None of those
inherited proofs uses the new target-directed rule.

\begin{lemma}[Elementary addition identities in $\mathcal T$]
\label{target:addition}
The calculus $\mathcal T$ proves
\[
 \begin{gathered}
 A(0,n)=n,\qquad A(Sa,n)=S A(a,n),\\
 A(a,0)=a,\qquad A(a,Sn)=S A(a,n),\qquad A(a,n)=A(n,a).
 \end{gathered}
\]
\end{lemma}
\begin{proof}
The first two are computation instances.  Induct on $a$ for $A(a,0)=a$,
then on $a$ for $A(a,Sn)=S A(a,n)$.  For commutativity, induct on $a$:
the base is $n=A(n,0)$ and the successor step is
$A(Sa,n)=S A(a,n)=S A(n,a)=A(n,Sa)$.
All inductions are on equations, using
Corollary~\ref{target:equation-induction}.
\end{proof}

\begin{theorem}[Identical consequence relations]
\label{target:exact-equivalence}
For every context $\Gamma$ and every formula $F$ in the full common
equation/$Q$/$D$ grammar,
\[
 \Gamma\vdash_{\mathcal T}F
 \quad\Longleftrightarrow\quad
 \Gamma\vdash_{\mathcal R}F.
\]
Both directions are effective transformations of finite derivations.
\end{theorem}
\begin{proof}
For $\mathcal T\to\mathcal R$, only a target-directed QE node needs to be
replaced.  If its matrix is an equation and its target is $Sr$, use the
inherited one-step rule at $r$.  If its matrix is an equation and its
target is not a successor, its minor is $E(0)$; use orbit QE at source $0$
and initializer $t$, then I0 and equation transport.  For a compound matrix,
its minor is again the zero instance; the same orbit specialization gives
$B(I(0,t;u.Su))$.  Use the already established full transport of
$\mathcal R$ and I0 to obtain $B(t)$.  The last step is why compound
transport must be placed before this comparison.

For $\mathcal R\to\mathcal T$, only orbit QE needs replacement.  For a
compound matrix, target-directed QE from $B(0)$ may target the exact term
$A(0,n)$, so it gives the required conclusion immediately.

For an equation matrix, retain $Q_yE(y)$ and $E(a)$.  Equation induction on
$k$ proves $E(A(k,a))$: the base follows from $A(0,a)=a$, and the successor
step follows from one-step elimination at $A(k,a)$ and
$A(Sk,a)=S A(k,a)$.  Specialize at $k=n$.  Commutativity from
Lemma~\ref{target:addition} and equation transport convert
$E(A(n,a))$ to $E(A(a,n))$, exactly the old target.
All other inference nodes translate literally.  The local constructions
retain arbitrary side assumptions and preserve fresh-eigenparameter
conditions after renaming.
\end{proof}

The addition identities in this last comparison explain how to simulate an
old orbit certificate.  They are not required by the new primitive rule,
the new checker, the induction bootstrap, or the direct proof of its
ordinary soundness.

\subsection{Ordinary interpretation and both full-language composites}
\label{target:biinterpretation-section}

The ordinary reading is exactly that of $\mathcal R$:
\[
 \begin{aligned}
 \llbracket s=t\rrbracket&=(s=t),\\
 \llbracket Q_yE\rrbracket
   &=\forall y\,(E(y)\to E(Sy))&&\text{if $E$ is an equation},\\
 \llbracket Q_yB\rrbracket
   &=\llbracket B(0)\rrbracket\to\forall y\,\llbracket B(y)\rrbracket
      &&\text{if $B$ is compound},\\
 \llbracket D_{x,y}B\rrbracket&=\exists x\,\llbracket Q_yB\rrbracket.
 \end{aligned}
\]
For completeness, the new QE is directly sound in the inherited scalar
ordinary theory $T_{\rm eq}$.  At an atomic successor target it uses one
progress implication.  At an atomic nonsuccessor target it uses equation
induction from the zero instance.  At a compound matrix it is ordinary
conditional-universal elimination.  Its soundness does not rely on the
target term denoting any particular kind of number.  The other rules have
the same soundness proofs as the inherited progressive calculus.

Here is the precise composition with the existing native/iPRA bridge.
Let $\mathcal N$ be the native calculus whose closures are uniformly read
as $B(0)\to\forall yB(y)$.  Define the inherited padding
\[
 G_v(s=t):=(I(v,0;u.s)=I(v,0;u.t)),\qquad
 L(E):=D_{r,v}G_vE,
\]
using fresh unused $u,r,v$ where required, and put $P(E)=L(E)$ for an
equation and $P(B)=B$ for a compound formula.  The inherited forward
translation $F:\mathcal N\to\mathcal R$ fixes equations and terms and
satisfies
\[
 F(Q_yB)=Q_yP(FB),\qquad F(D_{x,y}B)=D_{x,y}P(FB).
\]
The reverse translation $J:\mathcal R\to\mathcal N$ fixes equations and
terms.  At atomic closures it uses the native logical macros for
progressivity:
\[
 \begin{aligned}
 J(Q_yE)&=\forall_{\mathcal N}y\,
                    (E(y)\to_{\mathcal N}E(Sy)),\\
 J(D_{x,y}E)&=\exists_{\mathcal N}x\,\forall_{\mathcal N}y\,
                    (E(x,y)\to_{\mathcal N}E(x,Sy)).
 \end{aligned}
\]
At compound matrices it recurses without changing the binder:
$J(Q_yB)=Q_yJB$ and $J(D_{x,y}B)=D_{x,y}JB$.
The native macros and the full proofs of this inherited translation theorem
are retained in the source material; they are not new axioms of $\mathcal T$.

The inherited result\cite{uniform} proves derivation preservation in both directions and
both contextual formula composites,
\[
 A\dashv\vdash_{\mathcal N}JF(A),\qquad
 B\dashv\vdash_{\mathcal R}FJ(B),
\]
for every old $A$ and every progressive-language $B$.  In particular, its
second composite is a theorem for all unguarded atomic closures as well;
it is not restricted to a canonical public fragment.

\begin{theorem}[Whole-domain bi-interpretability with iPRA]
\label{target:biinterpretation}
The ordinary realization of $\mathcal T$ is equiconsistent and
bi-interpretable with full first-order intuitionistic primitive recursive
arithmetic.  Both numerical domains are whole, numerical equality is actual
equality, and both numerical composite isomorphisms are the identity graph
$x=y$.  Every formula of the unrestricted native grammar has the stated
round trip.
\end{theorem}
\begin{proof}
By Theorem~\ref{target:exact-equivalence}, the identity formula translation
between $\mathcal T$ and $\mathcal R$ preserves their entire consequence
relations.  Read the same syntactic forward map $F$ as taking values in
$\mathcal T$, and read the same reverse map $J$ as starting from
$\mathcal T$.  Translate the inherited $\mathcal R$ derivations by that
theorem.  This proves preservation and the composite
$B\dashv\vdash_{\mathcal T}FJ(B)$, while the $\mathcal N$ composite is
unchanged.

Compose these maps with the existing native/scalar/iPRA bi-interpretation.\cite{bridge}
The scalar terms, domain, and equality are unchanged by the present maps;
therefore the inherited numerical identity composites are unchanged.
The contradiction $0=S0$ is fixed, giving equiconsistency by the effective
proof translations.  This is an external theorem about the two formal
theories, not an assertion within iPRA of its own consistency.
\end{proof}

\paragraph{Scope of the result.}
The new mathematics is the target-directed primitive elimination,
its constant-size wrapper recovery of arbitrary-target closure, and the
exact full-language presentation theorem above.  The progressive
interpretation, padding translations, ordinary logical macros, and the
arithmetic bridge to iPRA are inherited and explicitly identified.
The proofs in this section are written metatheory, not a newly compiled
Lean verification.  In the matched standalone implementation, the new
checker improves the orbit checker by the recorded 21 normalized bytes,
11 significant tokens, and 8 AST nodes.  The already existing
witness-sensitive checker remains smaller in that comparison.  These facts
support a local refinement of the orbit presentation; they do not establish
a new minimum over every current proposal.
